% GENERATED FILE. Edit canonical lessons, then run npm run generate:books.
% canonical-source-hash: fnv1a64:b45a5cd0e78514a3
% canonical-lessons: MR-C359-vakil, MR-C359-khatla, MR-C359-sunavni, MR-C359-nivada, MR-C359-dava

\chapter{In Court}
\label{ch:mr-in-court}
\hlchaptermodality{359}

\begin{chapteropening}
\textbf{By the end of this chapter:} \emph{I can talk about lawyers, cases and verdicts.}

Complete the last lesson of chapter 359: I can talk about lawyers, cases and verdicts.
\end{chapteropening}

\section[\texorpdfstring{\mr{वकील} (vakīl)}{वकील (vakīl)}]{\mr{वकील} (vakīl) — a lawyer}

\label{lesson:MR-C359-vakil}



\begin{quote}
Before the new one: say the Marathi for to cancel, then the Marathi for to grant.
\end{quote}

\subsection*{You'll want to know: \mr{वकील}}
\textbf{\mr{वकील}} — \emph{vakīl} — \textquotedblleft{}a lawyer\textquotedblright{}.

\begin{quote}
\mr{घराच्या} \mr{वादासाठी} \mr{त्याने} \mr{वकील} \mr{नेमला}.
\end{quote}

\begin{quote}
\emph{gharācyā vādāsāṭhī tyāne vakīl nemlā.}
\end{quote}

\begin{quote}
He appointed a lawyer for the dispute over the house.
\end{quote}

Masculine.

Plural \mr{वकील}; \mr{वकिली}, the legal profession.

\begin{cousinweb}
\textbf{\mr{वकील}} (\emph{vakīl}) is a loan from Persian, usually traced to Arabic \emph{wakīl}, \textquotedblleft{}agent, one entrusted\textquotedblright{}.
\end{cousinweb}

\begin{grammarlens}[title={What you've built}]
One more word for giving reasons and opinions.
\end{grammarlens}

\subsection*{Your turn}
\begin{itemize}
\raggedright
  \item \emph{Say it:} \emph{vakīl}
  \item \emph{Say it:} \emph{vakīl}, once more
  \item \emph{Say it:} \emph{gharācyā vādāsāṭhī tyāne vakīl nemlā.} — He appointed a lawyer for the dispute over the house.
  \item \emph{Recall:} say the Marathi for to cancel, then the Marathi for to grant, then say \emph{vakīl} again
\end{itemize}

\begin{tcolorbox}[breakable,colback=teal!4,colframe=teal!35!black,title={Before you move on}]
What does \mr{वकील} mean? (\textquotedblleft{}A lawyer\textquotedblright{}.) Say it once more.
\end{tcolorbox}
\section[\texorpdfstring{\mr{खटला} (khaṭlā)}{खटला (khaṭlā)}]{\mr{खटला} (khaṭlā) — a lawsuit; a case}

\label{lesson:MR-C359-khatla}



\begin{quote}
Before the new one: say the Marathi for to grant, then the Marathi for a lawyer.
\end{quote}

\subsection*{You'll want to know: \mr{खटला}}
\textbf{\mr{खटला}} — \emph{khaṭlā} — \textquotedblleft{}a lawsuit; a case\textquotedblright{}.

\begin{quote}
\mr{हा} \mr{खटला} \mr{तीन} \mr{वर्षं} \mr{चालला}.
\end{quote}

\begin{quote}
\emph{hā khaṭlā tīn varṣaṁ cāllā.}
\end{quote}

\begin{quote}
This case went on for three years.
\end{quote}

Masculine. Plural: \mr{खटले}, \emph{khaṭle}.

\mr{खटला} \mr{चालणे}, a case is heard; \mr{खटला} \mr{भरणे}, to file a suit.

\begin{cousinweb}
\textbf{\mr{खटला}} (\emph{khaṭlā}) is a Marathi noun whose origin is not settled.
\end{cousinweb}

\begin{grammarlens}[title={What you've built}]
One more word for giving reasons and opinions.
\end{grammarlens}

\subsection*{Your turn}
\begin{itemize}
\raggedright
  \item \emph{Say it:} \emph{khaṭlā}
  \item \emph{Say it:} \emph{khaṭlā}, once more
  \item \emph{Say it:} \emph{hā khaṭlā tīn varṣaṁ cāllā.} — This case went on for three years.
  \item \emph{Recall:} say the Marathi for to grant, then the Marathi for a lawyer, then say \emph{khaṭlā} again
\end{itemize}

\begin{tcolorbox}[breakable,colback=teal!4,colframe=teal!35!black,title={Before you move on}]
What does \mr{खटला} mean? (\textquotedblleft{}A lawsuit; a case\textquotedblright{}.) Say it once more.
\end{tcolorbox}
\section[\texorpdfstring{\mr{सुनावणी} (sunāvṇī)}{सुनावणी (sunāvṇī)}]{\mr{सुनावणी} (sunāvṇī) — a hearing (in court)}

\label{lesson:MR-C359-sunavni}



\begin{quote}
Before the new one: say the Marathi for a lawyer, then the Marathi for a lawsuit.
\end{quote}

\subsection*{You'll want to know: \mr{सुनावणी}}
\textbf{\mr{सुनावणी}} — \emph{sunāvṇī} — \textquotedblleft{}a hearing (in court)\textquotedblright{}.

\begin{quote}
\mr{खटल्याची} \mr{पुढची} \mr{सुनावणी} \mr{सोमवारी} \mr{आहे}.
\end{quote}

\begin{quote}
\emph{khaṭalyācī puḍhcī sunāvṇī somvārī āhe.}
\end{quote}

\begin{quote}
The next hearing of the case is on Monday.
\end{quote}

Feminine.

Built on \mr{सुनावणे}, to pronounce: \mr{शिक्षा} \mr{सुनावणे}, to pass sentence.

\begin{cousinweb}
\textbf{\mr{सुनावणी}} (\emph{sunāvṇī}) is built on \mr{सुनावणे}, \textquotedblleft{}to tell to, to pronounce\textquotedblright{}, with the noun ending -\mr{णी}; the origin of the verb is not settled.
\end{cousinweb}

\begin{grammarlens}[title={What you've built}]
One more word for giving reasons and opinions.
\end{grammarlens}

\subsection*{Your turn}
\begin{itemize}
\raggedright
  \item \emph{Say it:} \emph{sunāvṇī}
  \item \emph{Say it:} \emph{sunāvṇī}, once more
  \item \emph{Say it:} \emph{khaṭalyācī puḍhcī sunāvṇī somvārī āhe.} — The next hearing of the case is on Monday.
  \item \emph{Recall:} say the Marathi for a lawyer, then the Marathi for a lawsuit, then say \emph{sunāvṇī} again
\end{itemize}

\begin{tcolorbox}[breakable,colback=teal!4,colframe=teal!35!black,title={Before you move on}]
What does \mr{सुनावणी} mean? (\textquotedblleft{}A hearing (in court)\textquotedblright{}.) Say it once more.
\end{tcolorbox}
\section[\texorpdfstring{\mr{निवाडा} (nivāḍā)}{निवाडा (nivāḍā)}]{\mr{निवाडा} (nivāḍā) — a verdict; a settlement}

\label{lesson:MR-C359-nivada}



\begin{quote}
Before the new one: say the Marathi for a lawsuit, then the Marathi for a hearing.
\end{quote}

\subsection*{You'll want to know: \mr{निवाडा}}
\textbf{\mr{निवाडा}} — \emph{nivāḍā} — \textquotedblleft{}a verdict; a settlement\textquotedblright{}.

\begin{quote}
\mr{न्यायाधीशांनी} \mr{शेवटी} \mr{निवाडा} \mr{दिला}.
\end{quote}

\begin{quote}
\emph{nyāyādhīśānnī śevṭī nivāḍā dilā.}
\end{quote}

\begin{quote}
The judge finally gave the verdict.
\end{quote}

Masculine.

\mr{निवाडा} \mr{देणे} or \mr{करणे}, to give a ruling.

\begin{cousinweb}
\textbf{\mr{निवाडा}} (\emph{nivāḍā}) is inherited: CDIAL lists Marathi \emph{nivāḍā}, \textquotedblleft{}settlement\textquotedblright{}, under \emph{*nirvarati}, \textquotedblleft{}restrains\textquotedblright{} (CDIAL 7392).
\end{cousinweb}

\begin{grammarlens}[title={What you've built}]
One more word for giving reasons and opinions.
\end{grammarlens}

\subsection*{Your turn}
\begin{itemize}
\raggedright
  \item \emph{Say it:} \emph{nivāḍā}
  \item \emph{Say it:} \emph{nivāḍā}, once more
  \item \emph{Say it:} \emph{nyāyādhīśānnī śevṭī nivāḍā dilā.} — The judge finally gave the verdict.
  \item \emph{Recall:} say the Marathi for a lawsuit, then the Marathi for a hearing, then say \emph{nivāḍā} again
\end{itemize}

\begin{tcolorbox}[breakable,colback=teal!4,colframe=teal!35!black,title={Before you move on}]
What does \mr{निवाडा} mean? (\textquotedblleft{}A verdict; a settlement\textquotedblright{}.) Say it once more.
\end{tcolorbox}
\section[\texorpdfstring{\mr{दावा} (dāvā)}{दावा (dāvā)}]{\mr{दावा} (dāvā) — a claim; a lawsuit}

\label{lesson:MR-C359-dava}



\begin{quote}
Before the new one: say the Marathi for a hearing, then the Marathi for a verdict.
\end{quote}

\subsection*{You'll want to know: \mr{दावा}}
\textbf{\mr{दावा}} — \emph{dāvā} — \textquotedblleft{}a claim; a lawsuit\textquotedblright{}.

\begin{quote}
\mr{भावाने} \mr{घरावर} \mr{दावा} \mr{केला}.
\end{quote}

\begin{quote}
\emph{bhāvāne gharāvar dāvā kelā.}
\end{quote}

\begin{quote}
The brother laid claim to the house.
\end{quote}

Masculine. Plural: \mr{दावे}, \emph{dāve}.

X \mr{वर} \mr{दावा} \mr{करणे}, to lay claim to X; \mr{दावा} \mr{दाखल} \mr{करणे}, to file a suit.

\begin{cousinweb}
\textbf{\mr{दावा}} (\emph{dāvā}) is a loan from Persian, usually traced to Arabic \emph{daʿwā}, \textquotedblleft{}a claim, a lawsuit\textquotedblright{}.
\end{cousinweb}

\begin{grammarlens}[title={What you've built}]
One more word for giving reasons and opinions.
\end{grammarlens}

\subsection*{Your turn}
\begin{itemize}
\raggedright
  \item \emph{Say it:} \emph{dāvā}
  \item \emph{Say it:} \emph{dāvā}, once more
  \item \emph{Say it:} \emph{bhāvāne gharāvar dāvā kelā.} — The brother laid claim to the house.
  \item \emph{Recall:} say the Marathi for a hearing, then the Marathi for a verdict, then say \emph{dāvā} again
\end{itemize}

\begin{tcolorbox}[breakable,colback=teal!4,colframe=teal!35!black,title={Before you move on}]
What does \mr{दावा} mean? (\textquotedblleft{}A claim; a lawsuit\textquotedblright{}.) Say it once more.
\end{tcolorbox}
